\documentclass[10pt,a4paper]{amsart}
\usepackage[margin=19mm]{geometry}
\usepackage{amsmath,amssymb,mathtools}
\usepackage[scr=rsfs,cal=euler]{mathalfa}
\usepackage{xfrac}
\usepackage[unicode,hidelinks,bookmarks=false]{hyperref}
\newcommand{\ie}{i.e.\ }
\newcommand{\cf}{cf.\ }
\newcommand{\resp}{respectively\ }
\newcommand{\Subsection}{\subsection}
\providecommand{\nobreakdash}{\nobreak\hbox{-}}
\setlength{\emergencystretch}{2em}
\multlinegap0pt
\usepackage{amssymb}
\usepackage{enumerate}
\usepackage[shortlabels]{enumitem}
\usepackage[normalem]{ulem}
\newtheorem{lemma}{Lemma}
\title{Clique factors in random samplings of regular graphs}
\author{Wanting Sun, Shunan Wei, Donglei Yang}
\date{Source: arXiv:2512.20287v1 (CC BY 4.0)}
\begin{document}
\maketitle

\noindent\emph{Source and licence.} Clique factors in random samplings of regular graphs. Version arXiv:2512.20287v1 (CC BY 4.0), distributed by arXiv under the Creative Commons Attribution 4.0 International licence, http://creativecommons.org/licenses/by/4.0/, as recorded in the arXiv licence metadata for this exact version and shown on the abstract page https://arxiv.org/abs/2512.20287v1. Full licence notice: \texttt{licenses/arxiv-2512.20287-CC-BY-4.0.txt}.\par\smallskip

\noindent\textit{Editorial scope.} Counted unit: the article's lemma vertex-cover (source0.tex lines 535--537). The article's complete original proof of it is delivered with it (source0.tex lines 1647--1698, one \texttt{proof} environment of 52 lines, which the article places away from the statement). No mathematics is written, repaired or re-derived: the statement and the proof are the article's own lines, in the article's own notation, and no retained line is substituted or re-flowed.  No further source-local context is needed: every cross-reference of every retained line resolves inside the two delivered spans.  Nothing was omitted from any retained span, so the package carries no omission record.  No retained line contains a citation, so the wrapper carries no bibliography and no bibliography entry.  No retained line contains a figure environment, an image-inclusion command or drawing code, so no image is delivered and the package declares no asset dependency.  Editorial formatting only, in addition: an AMS \texttt{amsart} preamble that restates the article's own declarations and macros used by the retained lines, namely the environments \texttt{lemma} on separate counters, together with the standard packages those lines need.  The retained environments print in this document as Lemma 1 in order of appearance, whereas the article prints the same items with the numbering of its own full document.  The complete native source of the article as distributed in the arXiv e-print for this exact version is bundled unchanged as \texttt{sources/191550/source0.tex}.
\begin{lemma}\label{vertex-cover}
Let $G$ be an $((r-1)n+1)$-regular graph on $rn$ vertices and $\{A_1,A_2,\ldots,A_r\}$ be a balanced partition of $V(G)$. Suppose that $C_i$ is a minimum vertex cover of $G[A_i]$ for each $i\in [r]$. Then $\max\{|C_i|:i\in [r]\}\geq \frac{1}{r}\sqrt{n}$.
\end{lemma}

\begin{proof}[\bf Proof of Lemma~\ref{vertex-cover}]
Let $G$ be an $((r-1)n+1)$-regular graph on $rn$ vertices, and let $\{A_1,A_2,\ldots,A_r\}$ be a balanced partition of $V(G)$. Suppose that $A_i^1$ is a minimum vertex cover of $G[A_i]$ for each $i\in [r]$. Without loss of generality, assume that $|A_1^1|\geq |A_2^1|\geq \cdots \geq |A_r^1|$. Clearly, there exists a vertex cover $C_i$ of $G[A_i]$ with size $|A_2^1|$ for each $i\in [2,r]$. For convenience, let $C_1:=A_1^1$. Denote $x_i:=|C_i|$ and $B_i:=A_i\setminus C_i$ for each $i\in [r]$. Hence, $G[B_i]=\emptyset$ for each $i\in [r]$. We proceed by considering the following two possible cases. 
% ~{(I suggest to update $C_i$ instead of $A_i^1$, and use $B_i$ for $A_i\setminus C_i$)}

\medskip
{\bf Case 1. $e\Big(B_1,\bigcup_{j\neq 2}C_j\Big)\leq e\Big(C_2,\bigcup_{j\neq 1}B_j\Big)$.}
\medskip

Since $G$ is  $((r-1)n+1)$-regular, one has 
$$
e(C_2,B_1)\leq |C_2|((r-1)n+1)-e\Big(C_2,\bigcup_{j\neq 1}B_j\Big).
$$
Note that as $x_2=x_3=\cdots=x_r$,
$$
e\Big(B_1,\bigcup_{j\neq 1}B_j\Big)\le\sum_{j\neq1}(n-x_1)(n-x_j)=(n-x_1)(r-1)(n-x_2).
$$
Therefore, we have
\begin{align*}
e\Big(B_1,\bigcup_{j\neq 2}C_j\Big)&=|B_1|((r-1)n+1)-e(B_1,C_2)-e\Big(B_1,\bigcup_{j\neq 1}B_j\Big)\\
&\geq (n-x_1)((r-1)n+1)-x_2((r-1)n+1)+e\Big(C_2,\bigcup_{j\neq 1}B_j\Big)-e\Big(B_1,\bigcup_{j\neq 1}B_j\Big)\\
&\geq (n-x_1-x_2)((r-1)n+1)+e\Big(C_2,\bigcup_{j\neq 1}B_j\Big)-(n-x_1)(r-1)(n-x_2).
\end{align*}
%where the last equality holds as $x_j=x_2$ for each $j\in [2,r]$.
By our assumption in this case, one has
\begin{align*}
(n-x_1-x_2)((r-1)n+1)\leq (n-x_1)(r-1)(n-x_2).
\end{align*}
It follows that $n\leq (r-1)x_1x_2+(x_1+x_2)\leq (r-1)x_1^2+2x_1.$ Thus, $x_1\geq \frac{1}{r}\sqrt{n}$, as desired.

\medskip
{\bf Case 2. $e\Big(B_1,\bigcup_{j\neq 2}C_j\Big)> e\Big(C_2,\bigcup_{j\neq 1}B_j\Big)$.}
\medskip

Notice that
$$
e\Big(\bigcup_{j\neq 2}C_j,\bigcup_{j\neq 1}B_j\Big)\leq \Big|\bigcup_{j\neq 2}C_j\Big|((r-1)n+1)-e\Big(\bigcup_{j\neq 2}C_j,B_1\Big).
$$
Therefore,\allowdisplaybreaks
\begin{align*}
e\Big(\bigcup_{j\neq 1}B_j,C_2\Big)=&\Big|\bigcup_{j\neq 1}B_j\Big|((r-1)n+1)-e\Big(\bigcup_{j\neq 1}B_j,\bigcup_{j\neq 2}C_j\Big)-e\Big(\bigcup_{j\neq 1}B_j,\bigcup_{t\in [r]}B_t\Big)\\
\geq &\Big|\bigcup_{j\neq 1}B_j\Big|((r-1)n+1)-\Big|\bigcup_{j\neq 2}C_j\Big|((r-1)n+1)+e\Big(\bigcup_{j\neq 2}C_j,B_1\Big)-\sum_{j\neq 1}\sum_{t\neq j}e(B_j,B_t)\\
\geq &((r-1)n+1)\Big((r-1)n-\sum_{j\neq 1}x_j-\sum_{j\neq 2}x_j\Big)+e\Big(\bigcup_{j\neq 2}C_j,B_1\Big)-\sum_{j\neq 1}\sum_{t\neq j}(n-x_j)(n-x_t)\\
=&((r-1)n+1)\Big((r-1)n-(r-1)x_2-x_1-(r-2)x_2\Big)+e\Big(\bigcup_{j\neq 2}C_j,B_1\Big)\\
&-(r-1)(n-x_2)\Big((r-1)n-x_1-(r-2)x_2\Big)\\
=&(r-1)n-x_1-(2r-3)x_2-(r-1)x_1x_2-(r-1)(r-2)x_2^2+e\Big(\bigcup_{j\neq 2}C_j,B_1\Big).
\end{align*}
By the assumption in this case, one has
$$
(r-1)n\leq x_1+(2r-3)x_2+(r-1)x_1x_2+(r-1)(r-2)x_2^2\leq (r-1)\big((r-1)x_1^2+2x_1\big).
$$
Thus, $x_1\geq \frac{1}{r}\sqrt{n}$, as desired. 
\end{proof}

\end{document}
